% --- Archon physics formalization source begin ---
% archon:physics
% archon:covers PhyXMiniProblems/problem_phyx_mini_0552.lean
% archon:source-report reports/phyx_mini/problem_phyx_mini_0552.source.json

\chapter{Physics problem phyx\_mini\_0552}
\label{ch:PhyXMiniProblems_problem_phyx_mini_0552}

\paragraph{Problem source.}
\#\# Physical scenario

An observer \textbackslash{}( O \textbackslash{}) is standing on a platform of length \textbackslash{}( D\_0 = 65 \textbackslash{}text\{ m\} \textbackslash{}) on a space station. A rocket passes at a relative speed of \textbackslash{}( 0.80c \textbackslash{}) moving parallel to the edge of the platform. The observer \textbackslash{}( O \textbackslash{}) notes that the front and back of the rocket simultaneously line up with the ends of the platform at a particular instant as shown in figure.

\#\# Question

According to \textbackslash{}( O' \textbackslash{}), how long does it take for observer \textbackslash{}( O \textbackslash{}) to pass the entire length of the rocket?

\#\# Answer choices

- A: A: \textbackslash{}( 0.36 \textbackslash{}mu\textbackslash{}text\{s\} \textbackslash{})

- B: B: \textbackslash{}( 0.27 \textbackslash{}mu\textbackslash{}text\{s\} \textbackslash{})

- C: C: \textbackslash{}( 0.54 \textbackslash{}mu\textbackslash{}text\{s\} \textbackslash{})

- D: D: \textbackslash{}( 0.45 \textbackslash{}mu\textbackslash{}text\{s\} \textbackslash{})

\#\# Figure caption (auxiliary; use the image as primary evidence)

The image depicts a diagram involving two reference frames, \textbackslash{}( O \textbackslash{}) and \textbackslash{}( O' \textbackslash{}).

- **Objects:**
  - An elongated object (depicted as an arrow or rod) in frame \textbackslash{}( O' \textbackslash{}), labeled with an overall length of 65 meters.
  - There is an arrow indicating motion to the left with a speed of \textbackslash{}( 0.8c \textbackslash{}) (where \textbackslash{}( c \textbackslash{}) is the speed of light).
  
- **Text:**
  - Above the object, there is the measurement "65 m."
  - The velocity is labeled as \textbackslash{}( 0.8c \textbackslash{}) with an arrow pointing to the left.
  
- **Labeling:**
  - The moving frame is marked as \textbackslash{}( O' \textbackslash{}).
  - The stationary reference frame is marked as \textbackslash{}( O \textbackslash{}).

- **Relationships:**
  - The object in frame \textbackslash{}( O' \textbackslash{}) is moving to the left relative to frame \textbackslash{}( O \textbackslash{}). 

This setup typically illustrates principles from special relativity, involving relative velocity and length contraction.

\#\# Dataset metadata

Relativity; Physical Model Grounding Reasoning, Multi-Formula Reasoning

\paragraph{Recorded answer/context.}
D: D: \textbackslash{}( 0.45 \textbackslash{}mu\textbackslash{}text\{s\} \textbackslash{})

\paragraph{Figure/image path.}
/root/proposal\_for\_physic/science-mango/phyx\_mini\_run/phyx\_data/test\_image/552.png

\paragraph{Formalization target.}
create a compiling Lean file with sorry bodies at `PhyXMiniProblems/problem\_phyx\_mini\_0552.lean`. The Lean declarations must preserve the physical quantities, dimensions or dimensional roles, figure labels, governing-law hypotheses, and final relation expressed by this problem.
Use Mathlib/Physlib names found through LeanExplore where available. If a physics API is missing, introduce faithful local abstractions rather than scalar placeholder aliases.

\begin{theorem}[Physics formalization target]
\label{thm:physics:phyx_mini_0552:target}
\lean{PhyXMiniProblems.ProblemPhyXMini0552.problem_phyx_mini_0552}
\uses{def:physics:phyx-mini-0552:phyxminiproblems-problemphyxmini0552-rocketplatformsetup, def:physics:phyx-mini-0552:phyxminiproblems-problemphyxmini0552-matchespassingrocketscenario, def:physics:phyx-mini-0552:phyxminiproblems-problemphyxmini0552-matchessuppliedrocketfigure, def:physics:phyx-mini-0552:phyxminiproblems-problemphyxmini0552-matchesproblemreadouts, def:physics:phyx-mini-0552:phyxminiproblems-problemphyxmini0552-hasphysicalrelativisticparameters, def:physics:phyx-mini-0552:phyxminiproblems-problemphyxmini0552-satisfiessimultaneousendpointalignment, def:physics:phyx-mini-0552:phyxminiproblems-problemphyxmini0552-satisfiesrocketframetraversalmeasurement, def:physics:phyx-mini-0552:phyxminiproblems-problemphyxmini0552-satisfiesrelativisticlengthandtraversallaws, def:physics:phyx-mini-0552:phyxminiproblems-problemphyxmini0552-matchesdisplayedtimechoice}
The assigned autoformalize agent should translate this physics problem into Lean declarations in the covered file, with theorem and lemma proof bodies written as `by sorry`.
\end{theorem}
\begin{proof}
This is an autoformalization task, not a proof task. Produce statements that can later be proved by the physics prover without weakening the physical model.
\end{proof}

% --- Archon declaration topology begin ---
\section{Declaration topology}
\begin{definition}[Length Quantity]
  \label{def:physics:phyx-mini-0552:phyxminiproblems-problemphyxmini0552-lengthquantity}
  \lean{PhyXMiniProblems.ProblemPhyXMini0552.LengthQuantity}
  A nonnegative, unit-independent physical length
\end{definition}
\begin{proof}
  This is the declared model definition of Length Quantity.
\end{proof}
\begin{definition}[Time Quantity]
  \label{def:physics:phyx-mini-0552:phyxminiproblems-problemphyxmini0552-timequantity}
  \lean{PhyXMiniProblems.ProblemPhyXMini0552.TimeQuantity}
  A nonnegative, unit-independent physical elapsed time
\end{definition}
\begin{proof}
  This is the declared model definition of Time Quantity.
\end{proof}
\begin{definition}[length Readout]
  \label{def:physics:phyx-mini-0552:phyxminiproblems-problemphyxmini0552-lengthreadout}
  \lean{PhyXMiniProblems.ProblemPhyXMini0552.lengthReadout}
  \uses{def:physics:phyx-mini-0552:phyxminiproblems-problemphyxmini0552-lengthquantity}
  Read a physical length in a selected length unit
\end{definition}
\begin{proof}
  This is the declared model definition of length Readout.
\end{proof}
\begin{definition}[length In Meters]
  \label{def:physics:phyx-mini-0552:phyxminiproblems-problemphyxmini0552-lengthinmeters}
  \lean{PhyXMiniProblems.ProblemPhyXMini0552.lengthInMeters}
  \uses{def:physics:phyx-mini-0552:phyxminiproblems-problemphyxmini0552-lengthquantity, def:physics:phyx-mini-0552:phyxminiproblems-problemphyxmini0552-lengthreadout}
  Meter readout used by the 65 m label
\end{definition}
\begin{proof}
  This is the declared model definition of length In Meters.
\end{proof}
\begin{definition}[time Readout]
  \label{def:physics:phyx-mini-0552:phyxminiproblems-problemphyxmini0552-timereadout}
  \lean{PhyXMiniProblems.ProblemPhyXMini0552.timeReadout}
  \uses{def:physics:phyx-mini-0552:phyxminiproblems-problemphyxmini0552-timequantity}
  Read a physical elapsed time in a selected time unit
\end{definition}
\begin{proof}
  This is the declared model definition of time Readout.
\end{proof}
\begin{definition}[time In Seconds]
  \label{def:physics:phyx-mini-0552:phyxminiproblems-problemphyxmini0552-timeinseconds}
  \lean{PhyXMiniProblems.ProblemPhyXMini0552.timeInSeconds}
  \uses{def:physics:phyx-mini-0552:phyxminiproblems-problemphyxmini0552-timequantity, def:physics:phyx-mini-0552:phyxminiproblems-problemphyxmini0552-timereadout}
  Second readout used for event-time coordinates
\end{definition}
\begin{proof}
  This is the declared model definition of time In Seconds.
\end{proof}
\begin{definition}[time In Microseconds]
  \label{def:physics:phyx-mini-0552:phyxminiproblems-problemphyxmini0552-timeinmicroseconds}
  \lean{PhyXMiniProblems.ProblemPhyXMini0552.timeInMicroseconds}
  \uses{def:physics:phyx-mini-0552:phyxminiproblems-problemphyxmini0552-timequantity, def:physics:phyx-mini-0552:phyxminiproblems-problemphyxmini0552-timereadout}
  Microsecond readout used by the answer choices
\end{definition}
\begin{proof}
  This is the declared model definition of time In Microseconds.
\end{proof}
\begin{definition}[speed In Meters Per Second]
  \label{def:physics:phyx-mini-0552:phyxminiproblems-problemphyxmini0552-speedinmeterspersecond}
  \lean{PhyXMiniProblems.ProblemPhyXMini0552.speedInMetersPerSecond}
  Read a physical speed magnitude in SI meters per second
\end{definition}
\begin{proof}
  This is the declared model definition of speed In Meters Per Second.
\end{proof}
\begin{definition}[vacuum Speed Of Light In Meters Per Second]
  \label{def:physics:phyx-mini-0552:phyxminiproblems-problemphyxmini0552-vacuumspeedoflightinmeterspersecond}
  \lean{PhyXMiniProblems.ProblemPhyXMini0552.vacuumSpeedOfLightInMetersPerSecond}
  Physlib's exact vacuum speed of light, read in meters per second
\end{definition}
\begin{proof}
  This is the declared model definition of vacuum Speed Of Light In Meters Per Second.
\end{proof}
\begin{definition}[Inertial Frame Label]
  \label{def:physics:phyx-mini-0552:phyxminiproblems-problemphyxmini0552-inertialframelabel}
  \lean{PhyXMiniProblems.ProblemPhyXMini0552.InertialFrameLabel}
  The two inertial-frame labels printed in the supplied image
\end{definition}
\begin{proof}
  This is the declared model definition of Inertial Frame Label.
\end{proof}
\begin{definition}[Rocket Endpoint]
  \label{def:physics:phyx-mini-0552:phyxminiproblems-problemphyxmini0552-rocketendpoint}
  \lean{PhyXMiniProblems.ProblemPhyXMini0552.RocketEndpoint}
  The front (nose) and back (tail) endpoints of the left-moving rocket
\end{definition}
\begin{proof}
  This is the declared model definition of Rocket Endpoint.
\end{proof}
\begin{definition}[Platform Endpoint]
  \label{def:physics:phyx-mini-0552:phyxminiproblems-problemphyxmini0552-platformendpoint}
  \lean{PhyXMiniProblems.ProblemPhyXMini0552.PlatformEndpoint}
  The two ends of the horizontal station platform
\end{definition}
\begin{proof}
  This is the declared model definition of Platform Endpoint.
\end{proof}
\begin{definition}[Axis Direction]
  \label{def:physics:phyx-mini-0552:phyxminiproblems-problemphyxmini0552-axisdirection}
  \lean{PhyXMiniProblems.ProblemPhyXMini0552.AxisDirection}
  Direction along the common rocket/platform axis
\end{definition}
\begin{proof}
  This is the declared model definition of Axis Direction.
\end{proof}
\begin{definition}[Axis Relation]
  \label{def:physics:phyx-mini-0552:phyxminiproblems-problemphyxmini0552-axisrelation}
  \lean{PhyXMiniProblems.ProblemPhyXMini0552.AxisRelation}
  Relation of the rocket's motion axis to the platform edge
\end{definition}
\begin{proof}
  This is the declared model definition of Axis Relation.
\end{proof}
\begin{definition}[Figure Feature]
  \label{def:physics:phyx-mini-0552:phyxminiproblems-problemphyxmini0552-figurefeature}
  \lean{PhyXMiniProblems.ProblemPhyXMini0552.FigureFeature}
  Named features visible in the primary raster 552.png
\end{definition}
\begin{proof}
  This is the declared model definition of Figure Feature.
\end{proof}
\begin{definition}[Rocket Platform Figure]
  \label{def:physics:phyx-mini-0552:phyxminiproblems-problemphyxmini0552-rocketplatformfigure}
  \lean{PhyXMiniProblems.ProblemPhyXMini0552.RocketPlatformFigure}
  \uses{def:physics:phyx-mini-0552:phyxminiproblems-problemphyxmini0552-inertialframelabel, def:physics:phyx-mini-0552:phyxminiproblems-problemphyxmini0552-rocketendpoint, def:physics:phyx-mini-0552:phyxminiproblems-problemphyxmini0552-axisdirection, def:physics:phyx-mini-0552:phyxminiproblems-problemphyxmini0552-figurefeature}
  Qualitative and printed evidence in the primary raster. The scalar entries are annotations only; the corresponding physical quantities are independent fields of RocketPlatformSetup below
\end{definition}
\begin{proof}
  This is the declared model definition of Rocket Platform Figure.
\end{proof}
\begin{definition}[Rocket Platform Setup]
  \label{def:physics:phyx-mini-0552:phyxminiproblems-problemphyxmini0552-rocketplatformsetup}
  \lean{PhyXMiniProblems.ProblemPhyXMini0552.RocketPlatformSetup}
  \uses{def:physics:phyx-mini-0552:phyxminiproblems-problemphyxmini0552-lengthquantity, def:physics:phyx-mini-0552:phyxminiproblems-problemphyxmini0552-timequantity, def:physics:phyx-mini-0552:phyxminiproblems-problemphyxmini0552-inertialframelabel, def:physics:phyx-mini-0552:phyxminiproblems-problemphyxmini0552-rocketendpoint, def:physics:phyx-mini-0552:phyxminiproblems-problemphyxmini0552-platformendpoint, def:physics:phyx-mini-0552:phyxminiproblems-problemphyxmini0552-axisdirection, def:physics:phyx-mini-0552:phyxminiproblems-problemphyxmini0552-axisrelation, def:physics:phyx-mini-0552:phyxminiproblems-problemphyxmini0552-rocketplatformfigure}
  Independent physical quantities and frame-coordinate data for the scenario. rocketLengthInPlatformFrame is the simultaneous endpoint separation in frame O; rocketProperLength is the rocket's rest length in frame O'. The traversal time is an independent observable, not an answer-choice value
\end{definition}
\begin{proof}
  This is the declared model definition of Rocket Platform Setup.
\end{proof}
\begin{definition}[speed Fraction Of Light]
  \label{def:physics:phyx-mini-0552:phyxminiproblems-problemphyxmini0552-speedfractionoflight}
  \lean{PhyXMiniProblems.ProblemPhyXMini0552.speedFractionOfLight}
  \uses{def:physics:phyx-mini-0552:phyxminiproblems-problemphyxmini0552-speedinmeterspersecond, def:physics:phyx-mini-0552:phyxminiproblems-problemphyxmini0552-vacuumspeedoflightinmeterspersecond, def:physics:phyx-mini-0552:phyxminiproblems-problemphyxmini0552-rocketplatformsetup}
  The dimensionless relative-speed parameter β = v/c
\end{definition}
\begin{proof}
  This is the declared model definition of speed Fraction Of Light.
\end{proof}
\begin{definition}[lorentz Factor]
  \label{def:physics:phyx-mini-0552:phyxminiproblems-problemphyxmini0552-lorentzfactor}
  \lean{PhyXMiniProblems.ProblemPhyXMini0552.lorentzFactor}
  \uses{def:physics:phyx-mini-0552:phyxminiproblems-problemphyxmini0552-rocketplatformsetup, def:physics:phyx-mini-0552:phyxminiproblems-problemphyxmini0552-speedfractionoflight}
  Physlib's Lorentz factor for the relative speed
\end{definition}
\begin{proof}
  This is the declared model definition of lorentz Factor.
\end{proof}
\begin{definition}[Matches Passing Rocket Scenario]
  \label{def:physics:phyx-mini-0552:phyxminiproblems-problemphyxmini0552-matchespassingrocketscenario}
  \lean{PhyXMiniProblems.ProblemPhyXMini0552.MatchesPassingRocketScenario}
  \uses{def:physics:phyx-mini-0552:phyxminiproblems-problemphyxmini0552-rocketplatformsetup}
  The frame assignments and qualitative relative motion in the prose
\end{definition}
\begin{proof}
  This is the declared model definition of Matches Passing Rocket Scenario.
\end{proof}
\begin{definition}[Matches Supplied Rocket Figure]
  \label{def:physics:phyx-mini-0552:phyxminiproblems-problemphyxmini0552-matchessuppliedrocketfigure}
  \lean{PhyXMiniProblems.ProblemPhyXMini0552.MatchesSuppliedRocketFigure}
  \uses{def:physics:phyx-mini-0552:phyxminiproblems-problemphyxmini0552-rocketplatformfigure}
  Direct primary-image evidence: the rocket above the platform, both frame labels, the 65 m bracket spanning nose to tail, and a leftward 0.8 c velocity arrow
\end{definition}
\begin{proof}
  This is the declared model definition of Matches Supplied Rocket Figure.
\end{proof}
\begin{definition}[Matches Problem Readouts]
  \label{def:physics:phyx-mini-0552:phyxminiproblems-problemphyxmini0552-matchesproblemreadouts}
  \lean{PhyXMiniProblems.ProblemPhyXMini0552.MatchesProblemReadouts}
  \uses{def:physics:phyx-mini-0552:phyxminiproblems-problemphyxmini0552-lengthinmeters, def:physics:phyx-mini-0552:phyxminiproblems-problemphyxmini0552-rocketplatformsetup, def:physics:phyx-mini-0552:phyxminiproblems-problemphyxmini0552-speedfractionoflight}
  The prose and image annotations are linked to the dimensionful setup. The rocket's platform-frame length is linked operationally through the endpoint alignment predicate below, rather than being defined from the desired time
\end{definition}
\begin{proof}
  This is the declared model definition of Matches Problem Readouts.
\end{proof}
\begin{definition}[Has Physical Relativistic Parameters]
  \label{def:physics:phyx-mini-0552:phyxminiproblems-problemphyxmini0552-hasphysicalrelativisticparameters}
  \lean{PhyXMiniProblems.ProblemPhyXMini0552.HasPhysicalRelativisticParameters}
  \uses{def:physics:phyx-mini-0552:phyxminiproblems-problemphyxmini0552-lengthinmeters, def:physics:phyx-mini-0552:phyxminiproblems-problemphyxmini0552-speedinmeterspersecond, def:physics:phyx-mini-0552:phyxminiproblems-problemphyxmini0552-rocketplatformsetup, def:physics:phyx-mini-0552:phyxminiproblems-problemphyxmini0552-speedfractionoflight}
  Positivity, endpoint ordering, and the physical subluminal regime
\end{definition}
\begin{proof}
  This is the declared model definition of Has Physical Relativistic Parameters.
\end{proof}
\begin{definition}[Satisfies Simultaneous Endpoint Alignment]
  \label{def:physics:phyx-mini-0552:phyxminiproblems-problemphyxmini0552-satisfiessimultaneousendpointalignment}
  \lean{PhyXMiniProblems.ProblemPhyXMini0552.SatisfiesSimultaneousEndpointAlignment}
  \uses{def:physics:phyx-mini-0552:phyxminiproblems-problemphyxmini0552-lengthinmeters, def:physics:phyx-mini-0552:phyxminiproblems-problemphyxmini0552-rocketplatformsetup}
  Operational form of the simultaneous alignment observed in frame O. The rocket's nose is at the platform's left end and its tail at the right end at the same O-frame time. Both measured lengths are coordinate separations
\end{definition}
\begin{proof}
  This is the declared model definition of Satisfies Simultaneous Endpoint Alignment.
\end{proof}
\begin{definition}[Satisfies Rocket Frame Traversal Measurement]
  \label{def:physics:phyx-mini-0552:phyxminiproblems-problemphyxmini0552-satisfiesrocketframetraversalmeasurement}
  \lean{PhyXMiniProblems.ProblemPhyXMini0552.SatisfiesRocketFrameTraversalMeasurement}
  \uses{def:physics:phyx-mini-0552:phyxminiproblems-problemphyxmini0552-timeinseconds, def:physics:phyx-mini-0552:phyxminiproblems-problemphyxmini0552-rocketplatformsetup}
  Operational meaning of the elapsed time asked for in frame O': it is the difference between the times when observer O passes the rocket's front and back endpoints
\end{definition}
\begin{proof}
  This is the declared model definition of Satisfies Rocket Frame Traversal Measurement.
\end{proof}
\begin{definition}[Satisfies Relativistic Length And Traversal Laws]
  \label{def:physics:phyx-mini-0552:phyxminiproblems-problemphyxmini0552-satisfiesrelativisticlengthandtraversallaws}
  \lean{PhyXMiniProblems.ProblemPhyXMini0552.SatisfiesRelativisticLengthAndTraversalLaws}
  \uses{def:physics:phyx-mini-0552:phyxminiproblems-problemphyxmini0552-lengthreadout, def:physics:phyx-mini-0552:phyxminiproblems-problemphyxmini0552-rocketplatformsetup, def:physics:phyx-mini-0552:phyxminiproblems-problemphyxmini0552-lorentzfactor}
  The two governing laws used by the solution: longitudinal length contraction, L = L₀ / γ(β); and constant-speed traversal in the rocket frame, t v = L₀. They are stated for arbitrary compatible unit choices and contain no numerical traversal time or answer-choice label
\end{definition}
\begin{proof}
  This is the declared model definition of Satisfies Relativistic Length And Traversal Laws.
\end{proof}
\begin{definition}[Answer Choice]
  \label{def:physics:phyx-mini-0552:phyxminiproblems-problemphyxmini0552-answerchoice}
  \lean{PhyXMiniProblems.ProblemPhyXMini0552.AnswerChoice}
  Labels of the four microsecond-valued choices in the source
\end{definition}
\begin{proof}
  This is the declared model definition of Answer Choice.
\end{proof}
\begin{definition}[displayed Time Microseconds]
  \label{def:physics:phyx-mini-0552:phyxminiproblems-problemphyxmini0552-displayedtimemicroseconds}
  \lean{PhyXMiniProblems.ProblemPhyXMini0552.displayedTimeMicroseconds}
  \uses{def:physics:phyx-mini-0552:phyxminiproblems-problemphyxmini0552-answerchoice}
  Displayed traversal time in microseconds, represented exactly
\end{definition}
\begin{proof}
  This is the declared model definition of displayed Time Microseconds.
\end{proof}
\begin{definition}[recorded Dataset Answer]
  \label{def:physics:phyx-mini-0552:phyxminiproblems-problemphyxmini0552-recordeddatasetanswer}
  \lean{PhyXMiniProblems.ProblemPhyXMini0552.recordedDatasetAnswer}
  \uses{def:physics:phyx-mini-0552:phyxminiproblems-problemphyxmini0552-answerchoice}
  The answer label recorded by the source dataset
\end{definition}
\begin{proof}
  This is the declared model definition of recorded Dataset Answer.
\end{proof}
\begin{definition}[Matches Displayed Time Choice]
  \label{def:physics:phyx-mini-0552:phyxminiproblems-problemphyxmini0552-matchesdisplayedtimechoice}
  \lean{PhyXMiniProblems.ProblemPhyXMini0552.MatchesDisplayedTimeChoice}
  \uses{def:physics:phyx-mini-0552:phyxminiproblems-problemphyxmini0552-timeinmicroseconds, def:physics:phyx-mini-0552:phyxminiproblems-problemphyxmini0552-rocketplatformsetup, def:physics:phyx-mini-0552:phyxminiproblems-problemphyxmini0552-answerchoice, def:physics:phyx-mini-0552:phyxminiproblems-problemphyxmini0552-displayedtimemicroseconds}
  Agreement with a two-decimal-place time display. The half-step tolerance 0.005 μs accounts for rounding; the laws with Physlib's exact value of c give approximately 0.4517 μs, not exactly 0.45 μs
\end{definition}
\begin{proof}
  This is the declared model definition of Matches Displayed Time Choice.
\end{proof}
% --- Archon declaration topology end ---
% --- Archon physics formalization source end ---
